\label{ch:sterile-theory}

This chapter develops the physics of the 3+1 sterile-neutrino extension,
from the short-baseline anomalies that originally motivated it to the exact
matter-potential structure a fourth, gauge-singlet mass eigenstate implies.
It is deliberately implementation-agnostic: \cref{ch:sterile-implementation}
maps every equation derived here onto the concrete \texttt{tpeanuts} code
that evaluates it.

\section{Motivation: short-baseline anomalies}

\begin{theorybox}
The three-flavour paradigm, with two independent mass-squared splittings
$\Delta m^2_{21}\sim7.4\times10^{-5}\ \mathrm{eV^2}$ and
$\Delta m^2_{31}\sim2.5\times10^{-3}\ \mathrm{eV^2}$, accounts for the vast
majority of solar, atmospheric, reactor, and accelerator oscillation data.
A handful of short-baseline anomalies, however, are not accommodated by
this picture and instead point to oscillations at a much larger
mass-squared splitting, $\Delta m^2\sim\mathcal{O}(1)\ \mathrm{eV^2}$,
requiring a third, independent splitting and hence (at least) a fourth mass
eigenstate:
\begin{itemize}
  \item \textbf{LSND}: excess $\bar\nu_\mu\to\bar\nu_e$ appearance at
    $L/E\sim1$~m/MeV~\parencite{Aguilar2001LSND}.
  \item \textbf{MiniBooNE}: an independent, statistically significant
    excess of electron-like events in both neutrino and antineutrino mode,
    broadly consistent with the LSND
    signal~\parencite{AguilarArevalo2018MiniBooNE}.
  \item \textbf{The reactor antineutrino anomaly}: a $\sim6\%$ deficit
    between measured and predicted $\bar\nu_e$ flux at short baseline from
    reactor experiments, once reactor flux predictions were
    revised~\parencite{Mention2011ReactorAnomaly}.
  \item \textbf{The gallium anomaly}: a deficit in the $\nu_e$ event rate
    measured during the radioactive-source calibration campaigns of the
    GALLEX and SAGE solar-neutrino
    experiments~\parencite{GiuntiLaveder2011GalliumAnomaly}.
\end{itemize}
The simplest new-physics explanation common to all four is a single
additional mass eigenstate $\nu_4$, mostly composed of a new flavour
eigenstate $\nu_s$ that does not participate in the weak interaction at
all (a ``sterile'' neutrino) -- the ``3+1'' scheme. The appearance
anomalies (LSND, MiniBooNE) and the disappearance anomalies (reactor,
gallium) probe different combinations of the new mixing angles, and
existing global fits are in tension between the two classes; this tension,
and the different individual mixing angles preferred by different
data sets, is exactly why \cref{ch:sterile-observables} offers several
distinct named presets rather than a single best fit.
\end{theorybox}

\section{The 3+1 mixing matrix}

\begin{theorybox}
A fourth mass eigenstate $\nu_4$ and a fourth sterile flavour eigenstate
$\nu_s$ are added to the three active ones. Because $\nu_s$ is a
gauge-singlet, it does not receive the charged-current matter potential of
\S\ref{sec:sterile-matter} below, and the $4\times4$ mixing matrix is
parametrised, following the convention of
\textcite{KoppMachadoMaltoniSchwetz2013}, by adding three active--sterile
rotations $R_{14},R_{24},R_{34}$ to the standard 3-flavour structure,
\begin{equation}
  \keyresult{U_4 = R_{23}\,\Delta_4\,R_{24}\,R_{34}\,R_{13}\,\Delta_4^\dagger\,R_{12}\,R_{14}},
  \label{eq:u4}
\end{equation}
where $\theta_{14},\theta_{24},\theta_{34}$ are the new active--sterile
mixing angles and $\delta_{14},\delta_{24}$ their associated CP phases.
\Cref{eq:u4} reduces exactly to the 3-flavour Standard Model matrix
(\emph{embedded} in the larger $4\times4$ space) when
$\theta_{14}=\theta_{24}=\theta_{34}=0$: the sterile row/column decouple
completely and every propagation result of the pure Standard Model is
recovered without approximation.
\end{theorybox}

\begin{notebox}
\textbf{CP-phase counting.} For $N$ Dirac generations there are
$(N-1)(N-2)/2$ independent physical CP phases. For $N=4$ this gives exactly
\keyresult{three}: the standard phase $\delta_{13}$ (carried through
$R_{13}$ by $\Delta_4$, exactly as in the 3-flavour case), the $e$--$s$
phase $\delta_{14}$ (carried by $R_{14}$), and the $\mu$--$s$ phase
$\delta_{24}$ (carried by $R_{24}$). $R_{34}$ is always a \emph{real}
rotation: the fourth possible phase can be absorbed into a redefinition of
the charged-lepton fields, so no sterile preset (\cref{ch:sterile-observables})
ever needs a $\delta_{34}$ value.
\end{notebox}

\section{The matter potential: charged current versus neutral current}
\label{sec:sterile-matter}

\begin{theorybox}
Because $\nu_s$ is a gauge singlet, it never receives the charged-current
potential: in the extended flavour basis $(e,\mu,\tau,s)$, the CC term is
simply the 3-flavour Wolfenstein potential~\parencite{Wolfenstein1978},
zero-padded,
\begin{equation}
  H_{\rm mat}^{\rm CC} = \diag(V_{CC},0,0,0), \qquad V_{CC}=\pm\sqrt2\,G_F\,n_e\,L_{\rm scale}.
\end{equation}
The three \emph{active} flavours, in contrast, all share the same
neutral-current potential $V_{NC}=\mp\frac{1}{\sqrt2}G_F\,n_n\,L_{\rm scale}$
equally (lepton universality of the NC coupling). In the pure 3-flavour
Standard Model this shared term is an unobservable global phase and is
correctly never computed at all. Once a sterile state is present -- an
NC-blind fourth entry -- it is no longer a pure phase: writing out the full
diagonal and subtracting the common piece,
\begin{equation}
  \keyresult{
    \diag(V_{CC}+V_{NC},\,V_{NC},\,V_{NC},\,0) - V_{NC}\,I_4
    = \diag(V_{CC},\,0,\,0,\,-V_{NC})
  },
  \label{eq:sterile-nc}
\end{equation}
a genuine, physically observable relative potential of magnitude comparable
to $V_{CC}$ itself survives on the sterile diagonal entry. This term is
\emph{optional} in the sense that omitting it (the default) recovers the
CC-only sterile Hamiltonian used throughout most of the sterile-oscillation
literature and throughout earlier phases of this project -- unlike the
purely 3-flavour case, though, omitting it is a genuine approximation, not
an exact simplification, once $\theta_{24}$ or $\theta_{34}$ (which mix the
sterile index into the $\mu$/$\tau$ sector via the outer block,
\cref{ch:sterile-implementation}) are non-zero.
\end{theorybox}

\begin{notebox}
This sterile neutral-current term is structurally independent of, but can
be combined freely with, Non-Standard Interactions (a separate BSM
extension; see the companion \texttt{BSM\_NSI} document): both consume the
same underlying local neutron density $n_n(x)$, but they are additive,
independently switchable contributions to the matter Hamiltonian.
\end{notebox}

\section{The two-flavour disappearance limit}

\begin{theorybox}
When every Standard Model active mixing angle is switched off
($\theta_{12}=\theta_{13}=\theta_{23}=0$), the $\nu_e$ survival probability
reduces exactly to the textbook two-flavour vacuum formula
\begin{equation}
  \keyresult{P_{ee} = 1 - \sin^2(2\theta_{14})\,\sin^2\!\left(1.267\,\frac{\Delta m^2_{41}\,[\mathrm{eV^2}]\,L\,[\mathrm{km}]}{E\,[\mathrm{GeV}]}\right)},
\end{equation}
for \emph{any} value of $\theta_{24}$, since $\theta_{24}$ lives entirely
in the outer block $O_4$ (\cref{ch:sterile-implementation}), which leaves
the electron index untouched. The analogous statement for $\nu_\mu$
disappearance, $|U_{\mu4}|^2=\sin^2\theta_{24}$, holds exactly only when
\emph{both} $\theta_{23}=0$ \emph{and} $\theta_{14}=0$: $R_{14}$ and
$R_{24}$ both act on the sterile index and do not commute in general, so a
small residual correction survives in the $\nu_\mu$ channel whenever
$\theta_{14}$ and $\theta_{24}$ are simultaneously non-zero -- a property
of the physics, not of the numerical implementation, and the subject of
\cref{sec:results-kinematics}'s quantitative study.
\end{theorybox}

\begin{refbox}
Short-baseline anomalies: \cite{Aguilar2001LSND,AguilarArevalo2018MiniBooNE,Mention2011ReactorAnomaly,GiuntiLaveder2011GalliumAnomaly}.
3+1 mixing parametrisation: \cite{KoppMachadoMaltoniSchwetz2013}. Global 3+1
fit: \cite{GariazzoGiuntiLavederLi2017}. IceCube atmospheric $\nu_\mu$
disappearance search: \cite{IceCube2020SterileNumuDisappearance}. Standard
matter potential: \cite{Wolfenstein1978}.
\end{refbox}
